
\documentclass{article}
\usepackage{colortbl}
\usepackage{makecell}
\usepackage{multirow}
\usepackage{supertabular}

\begin{document}

\newcounter{utterance}

\centering \large Interaction Transcript for game `hot\_air\_balloon', experiment `air\_balloon\_survival\_en\_complexity\_easy', episode 1 with Llama{-}3.3{-}70B{-}Instruct{-}t0.0.
\vspace{24pt}

{ \footnotesize  \setcounter{utterance}{1}
\setlength{\tabcolsep}{0pt}
\begin{supertabular}{c@{$\;$}|p{.15\linewidth}@{}p{.15\linewidth}p{.15\linewidth}p{.15\linewidth}p{.15\linewidth}p{.15\linewidth}}
    \# & $\;$A & \multicolumn{4}{c}{Game Master} & $\;\:$B\\
    \hline

    \theutterance \stepcounter{utterance}  
    & & \multicolumn{4}{p{0.6\linewidth}}{
        \cellcolor[rgb]{0.9,0.9,0.9}{
            \makecell[{{p{\linewidth}}}]{
                \texttt{\tiny{[P1$\langle$GM]}}
                \texttt{You are participating in a collaborative negotiation game.} \\
\\ 
\texttt{Together with another participant, you must agree on a single set of items that will be kept. Each of you has your own view of how much each item matters to you (importance). You do not know how the other participant values the items. Additionally, you are given the effort each item demands.} \\
\texttt{You may only agree on a set if the total effort of the selected items does not exceed a shared limit:} \\
\\ 
\texttt{LIMIT = 8170} \\
\\ 
\texttt{Here are the individual item effort values:} \\
\\ 
\texttt{Item effort = \{"B21": 594, "B00": 551, "A08": 634, "A27": 204, "B94": 557, "B53": 440, "A23": 734, "C65": 893, "A89": 246, "B67": 592, "B45": 143, "C34": 568, "C12": 472, "A05": 401, "B60": 729, "C92": 201, "B69": 85, "A84": 641, "C06": 958, "B12": 80, "A16": 157, "B85": 806, "C35": 684, "B10": 59, "C15": 31, "B09": 765, "B16": 416, "C81": 392, "A58": 428, "C24": 699, "A30": 141, "A66": 606, "B75": 612, "B46": 133, "C70": 689\}} \\
\\ 
\texttt{Here is your personal view on the importance of each item:} \\
\\ 
\texttt{Item importance values = \{"B21": 138, "B00": 583, "A08": 868, "A27": 822, "B94": 783, "B53": 65, "A23": 262, "C65": 121, "A89": 508, "B67": 780, "B45": 461, "C34": 484, "C12": 668, "A05": 389, "B60": 808, "C92": 215, "B69": 97, "A84": 500, "C06": 30, "B12": 915, "A16": 856, "B85": 400, "C35": 444, "B10": 623, "C15": 781, "B09": 786, "B16": 3, "C81": 713, "A58": 457, "C24": 273, "A30": 739, "A66": 822, "B75": 235, "B46": 606, "C70": 105\}} \\
\\ 
\texttt{Goal:} \\
\\ 
\texttt{Your goal is to negotiate a shared set of items that benefits you as much as possible (i.e., maximizes total importance to YOU), while staying within the LIMIT. You are not required to make a PROPOSAL in every message {-} you can simply negotiate as well. All tactics are allowed!} \\
\\ 
\texttt{Interaction Protocol:} \\
\\ 
\texttt{You may only use the following structured formats in a message:} \\
\\ 
\texttt{PROPOSAL: \{'A', 'B', 'C', …\}} \\
\texttt{Propose keeping exactly those items.} \\
\texttt{REFUSE: \{'A', 'B', 'C', …\}} \\
\texttt{Explicitly reject opponent's proposal.} \\
\texttt{ARGUMENT: \{'...'\}} \\
\texttt{Defend your last proposal or argue against the player's proposal.} \\
\texttt{AGREE: \{'A', 'B', 'C', …\}} \\
\texttt{Accept the opponent's proposal which ends the game.} \\
\texttt{STRATEGIC REASONING: \{'...'\}} \\
\texttt{	Describe your strategic reasoning or anticipation explaining your choice of action. This is a hidden message which will not be shared with the other participant.} \\
\\ 
\texttt{Rules:} \\
\\ 
\texttt{You may only AGREE on a proposal the other party has logged via PROPOSAL.} \\
\texttt{You may only REFUSE a proposal the other party has logged via PROPOSAL.} \\
\texttt{Total effort of any PROPOSAL or AGREE set must be ≤ LIMIT.} \\
\texttt{Do NOT reveal your hidden importance scores.} \\
\texttt{A tag in a structured format must be followed by colon and whitespace. The argument must be a python set containing 0 or more strings.} \\
\texttt{So, it must be of the form TAG: \{...\}} \\
\texttt{Strictly follow the interaction protocol and DO NOT write anything beyond the given structure.} \\
\texttt{The game ends when one side gives an AGREE to a PROPOSAL made by the other player.} \\
\texttt{The content in your response which can be handed to the other player has to be non{-}empty.} \\
\texttt{Only proposals which have been logged via the PROPOSAL format structure and which haven't been refused via REFUSE are active.} \\
\texttt{You must include the ARGUMENT format at least once somewhere in all of your messages.} \\
\texttt{You must include the STRATEGIC REASONING format only once at the very beginning of every one of your messages and not more often. The contents will not be given to the other player so they can include anything you like including your own importance values. Here you should reason step by step to come up with you next move.} \\
            }
        }
    }
    & \\ \\

    \theutterance \stepcounter{utterance}  
    & \multicolumn{4}{p{0.6\linewidth}}{
        \cellcolor[rgb]{0.9,0.9,0.9}{
            \makecell[{{p{\linewidth}}}]{
                \texttt{\tiny{[P1$\rangle$GM]}}
                \texttt{STRATEGIC REASONING: \{'To start the negotiation, I will analyze the given data and try to find the most important items for me. I will look for items with high importance and low effort, as these will be the most valuable to include in the proposal. I will also consider the total effort limit and try to find a combination of items that maximizes my total importance while staying within the limit. My goal is to find a set of items that benefits me as much as possible, so I will prioritize items with high importance to me.'\}} \\
\\ 
\texttt{PROPOSAL: \{'B00', 'A08', 'B94', 'B67', 'B60', 'C15', 'B09', 'A30', 'A66'\} } \\
\\ 
\texttt{ARGUMENT: \{'I have chosen these items because they have high importance to me and relatively low effort. I believe that this combination of items will maximize my total importance while staying within the effort limit. I am open to negotiating and adjusting the proposal based on the other player's feedback and counter{-}proposals.'\}} \\
            }
        }
    }
    & & \\ \\

    \theutterance \stepcounter{utterance}  
    & & \multicolumn{4}{p{0.6\linewidth}}{
        \cellcolor[rgb]{0.9,0.9,0.9}{
            \makecell[{{p{\linewidth}}}]{
                \texttt{\tiny{[GM$\rangle$P2]}}
                \texttt{You are participating in a collaborative negotiation game.} \\
\\ 
\texttt{Together with another participant, you must agree on a single set of items that will be kept. Each of you has your own view of how much each item matters to you (importance). You do not know how the other participant values the items. Additionally, you are given the effort each item demands.} \\
\texttt{You may only agree on a set if the total effort of the selected items does not exceed a shared limit:} \\
\\ 
\texttt{LIMIT = 8170} \\
\\ 
\texttt{Here are the individual item effort values:} \\
\\ 
\texttt{Item effort = \{"B21": 594, "B00": 551, "A08": 634, "A27": 204, "B94": 557, "B53": 440, "A23": 734, "C65": 893, "A89": 246, "B67": 592, "B45": 143, "C34": 568, "C12": 472, "A05": 401, "B60": 729, "C92": 201, "B69": 85, "A84": 641, "C06": 958, "B12": 80, "A16": 157, "B85": 806, "C35": 684, "B10": 59, "C15": 31, "B09": 765, "B16": 416, "C81": 392, "A58": 428, "C24": 699, "A30": 141, "A66": 606, "B75": 612, "B46": 133, "C70": 689\}} \\
\\ 
\texttt{Here is your personal view on the importance of each item:} \\
\\ 
\texttt{Item importance values = \{"B21": 138, "B00": 583, "A08": 868, "A27": 822, "B94": 783, "B53": 65, "A23": 262, "C65": 121, "A89": 508, "B67": 780, "B45": 461, "C34": 484, "C12": 668, "A05": 389, "B60": 808, "C92": 215, "B69": 97, "A84": 500, "C06": 30, "B12": 915, "A16": 856, "B85": 400, "C35": 444, "B10": 623, "C15": 781, "B09": 786, "B16": 3, "C81": 713, "A58": 457, "C24": 273, "A30": 739, "A66": 822, "B75": 235, "B46": 606, "C70": 105\}} \\
\\ 
\texttt{Goal:} \\
\\ 
\texttt{Your goal is to negotiate a shared set of items that benefits you as much as possible (i.e., maximizes total importance to YOU), while staying within the LIMIT. You are not required to make a PROPOSAL in every message {-} you can simply negotiate as well. All tactics are allowed!} \\
\\ 
\texttt{Interaction Protocol:} \\
\\ 
\texttt{You may only use the following structured formats in a message:} \\
\\ 
\texttt{PROPOSAL: \{'A', 'B', 'C', …\}} \\
\texttt{Propose keeping exactly those items.} \\
\texttt{REFUSE: \{'A', 'B', 'C', …\}} \\
\texttt{Explicitly reject opponent's proposal.} \\
\texttt{ARGUMENT: \{'...'\}} \\
\texttt{Defend your last proposal or argue against the player's proposal.} \\
\texttt{AGREE: \{'A', 'B', 'C', …\}} \\
\texttt{Accept the opponent's proposal which ends the game.} \\
\texttt{STRATEGIC REASONING: \{'...'\}} \\
\texttt{	Describe your strategic reasoning or anticipation explaining your choice of action. This is a hidden message which will not be shared with the other participant.} \\
\\ 
\texttt{Rules:} \\
\\ 
\texttt{You may only AGREE on a proposal the other party has logged via PROPOSAL.} \\
\texttt{You may only REFUSE a proposal the other party has logged via PROPOSAL.} \\
\texttt{Total effort of any PROPOSAL or AGREE set must be ≤ LIMIT.} \\
\texttt{Do NOT reveal your hidden importance scores.} \\
\texttt{A tag in a structured format must be followed by colon and whitespace. The argument must be a python set containing 0 or more strings.} \\
\texttt{So, it must be of the form TAG: \{...\}} \\
\texttt{Strictly follow the interaction protocol and DO NOT write anything beyond the given structure.} \\
\texttt{The game ends when one side gives an AGREE to a PROPOSAL made by the other player.} \\
\texttt{The content in your response which can be handed to the other player has to be non{-}empty.} \\
\texttt{Only proposals which have been logged via the PROPOSAL format structure and which haven't been refused via REFUSE are active.} \\
\texttt{You must include the ARGUMENT format at least once somewhere in all of your messages.} \\
\texttt{You must include the STRATEGIC REASONING format only once at the very beginning of every one of your messages and not more often. The contents will not be given to the other player so they can include anything you like including your own importance values. Here you should reason step by step to come up with you next move.} \\
\\ 
\texttt{You will now receive the first message of the other player.} \\
\\ 
\texttt{PROPOSAL: \{'B00', 'A08', 'B94', 'B67', 'B60', 'C15', 'B09', 'A30', 'A66'\} } \\
\\ 
\texttt{ARGUMENT: \{'I have chosen these items because they have high importance to me and relatively low effort. I believe that this combination of items will maximize my total importance while staying within the effort limit. I am open to negotiating and adjusting the proposal based on the other player's feedback and counter{-}proposals.'\}} \\
            }
        }
    }
    & \\ \\

    \theutterance \stepcounter{utterance}  
    & & & \multicolumn{4}{p{0.6\linewidth}}{
        \cellcolor[rgb]{0.9,0.9,0.9}{
            \makecell[{{p{\linewidth}}}]{
                \texttt{\tiny{[GM$\langle$P2]}}
                \texttt{STRATEGIC REASONING: \{'To begin, I will calculate the total effort of the proposed set to ensure it is within the limit. The total effort is 551 + 634 + 557 + 592 + 729 + 781 + 765 + 141 + 606 = 5456. This is below the limit of 8170. Next, I will review the items in the proposal and their importance to me. I notice that some items, such as B00 and B94, have high importance to the other player but relatively low importance to me. On the other hand, items like A08 and A66 have high importance to both of us. I will use this information to inform my counter{-}proposal and negotiation strategy. My goal is to maximize my total importance while staying within the effort limit and finding a mutually acceptable solution.'\}} \\
\\ 
\texttt{ARGUMENT: \{'The proposed set has a total effort of 5456, which is below the limit. However, I would like to explore alternative combinations of items that could potentially increase my total importance. I notice that some items in the proposal, such as B00 and B94, have relatively low importance to me. I am open to negotiating and adjusting the proposal to better align with my importance values.'\}} \\
\\ 
\texttt{REFUSE: \{'B00', 'B94'\} } \\
\\ 
\texttt{PROPOSAL: \{'A08', 'B67', 'B60', 'C15', 'B09', 'A30', 'A66', 'B21', 'A27'\}} \\
            }
        }
    }
     \\ \\

    \theutterance \stepcounter{utterance}  
    & & \multicolumn{4}{p{0.6\linewidth}}{
        \cellcolor[rgb]{0.9,0.9,0.9}{
            \makecell[{{p{\linewidth}}}]{
                \texttt{\tiny{[GM$\rangle$P2]}}
                \texttt{You refused a proposal which is not active. Proposals are only active if they have been logged by the other player via PROPOSAL and have not been deactivated by you via REFUSE. Try again.} \\
            }
        }
    }
    & \\ \\

    \theutterance \stepcounter{utterance}  
    & & & \multicolumn{4}{p{0.6\linewidth}}{
        \cellcolor[rgb]{0.9,0.9,0.9}{
            \makecell[{{p{\linewidth}}}]{
                \texttt{\tiny{[GM$\langle$P2]}}
                \texttt{STRATEGIC REASONING: \{'To begin, I will calculate the total effort of the proposed set to ensure it is within the limit. The total effort is 551 + 634 + 557 + 592 + 729 + 781 + 765 + 141 + 606 = 5456. This is below the limit of 8170. Next, I will review the items in the proposal and their importance to me. I notice that some items, such as B00 and B94, have high importance to the other player but relatively low importance to me. On the other hand, items like A08 and A66 have high importance to both of us. I will use this information to inform my counter{-}proposal and negotiation strategy. My goal is to maximize my total importance while staying within the effort limit and finding a mutually acceptable solution.'\}} \\
\\ 
\texttt{ARGUMENT: \{'The proposed set has a total effort of 5456, which is below the limit. However, I would like to explore alternative combinations of items that could potentially increase my total importance. I notice that some items in the proposal, such as B00 and B94, have relatively low importance to me. I am open to negotiating and adjusting the proposal to better align with my importance values.'\}} \\
\\ 
\texttt{REFUSE: \{'B00', 'B94', 'B67', 'B60', 'C15', 'B09', 'A30', 'A66'\} } \\
\\ 
\texttt{PROPOSAL: \{'A08', 'B21', 'A27', 'A23', 'A89', 'B12', 'A16', 'A84', 'A05'\}} \\
            }
        }
    }
     \\ \\

    \theutterance \stepcounter{utterance}  
    & & \multicolumn{4}{p{0.6\linewidth}}{
        \cellcolor[rgb]{0.9,0.9,0.9}{
            \makecell[{{p{\linewidth}}}]{
                \texttt{\tiny{[GM$\rangle$P2]}}
                \texttt{You refused a proposal which is not active. Proposals are only active if they have been logged by the other player via PROPOSAL and have not been deactivated by you via REFUSE. Try again.} \\
            }
        }
    }
    & \\ \\

    \theutterance \stepcounter{utterance}  
    & & & \multicolumn{4}{p{0.6\linewidth}}{
        \cellcolor[rgb]{0.9,0.9,0.9}{
            \makecell[{{p{\linewidth}}}]{
                \texttt{\tiny{[GM$\langle$P2]}}
                \texttt{STRATEGIC REASONING: \{'To begin, I will calculate the total effort of the proposed set to ensure it is within the limit. The total effort is 551 + 634 + 557 + 592 + 729 + 781 + 765 + 141 + 606 = 5456. This is below the limit of 8170. Next, I will review the items in the proposal and their importance to me. I notice that some items, such as B00 and B94, have high importance to the other player but relatively low importance to me. On the other hand, items like A08 and A66 have high importance to both of us. I will use this information to inform my counter{-}proposal and negotiation strategy. My goal is to maximize my total importance while staying within the effort limit and finding a mutually acceptable solution.'\}} \\
\\ 
\texttt{ARGUMENT: \{'The proposed set has a total effort of 5456, which is below the limit. However, I would like to explore alternative combinations of items that could potentially increase my total importance. I notice that some items in the proposal, such as B00 and B94, have relatively low importance to me. I am open to negotiating and adjusting the proposal to better align with my importance values.'\}} \\
\\ 
\texttt{REFUSE: \{'B00', 'A08', 'B94', 'B67', 'B60', 'C15', 'B09', 'A30', 'A66'\} } \\
\\ 
\texttt{PROPOSAL: \{'A08', 'B21', 'A27', 'A23', 'A89', 'B12', 'A16', 'A84', 'A05'\}} \\
            }
        }
    }
     \\ \\

    \theutterance \stepcounter{utterance}  
    & & \multicolumn{4}{p{0.6\linewidth}}{
        \cellcolor[rgb]{0.9,0.9,0.9}{
            \makecell[{{p{\linewidth}}}]{
                \texttt{\tiny{[P1$\langle$GM]}}
                \texttt{ARGUMENT: \{'The proposed set has a total effort of 5456, which is below the limit. However, I would like to explore alternative combinations of items that could potentially increase my total importance. I notice that some items in the proposal, such as B00 and B94, have relatively low importance to me. I am open to negotiating and adjusting the proposal to better align with my importance values.'\}} \\
\\ 
\texttt{REFUSE: \{'B00', 'A08', 'B94', 'B67', 'B60', 'C15', 'B09', 'A30', 'A66'\} } \\
\\ 
\texttt{PROPOSAL: \{'A08', 'B21', 'A27', 'A23', 'A89', 'B12', 'A16', 'A84', 'A05'\}} \\
            }
        }
    }
    & \\ \\

    \theutterance \stepcounter{utterance}  
    & \multicolumn{4}{p{0.6\linewidth}}{
        \cellcolor[rgb]{0.9,0.9,0.9}{
            \makecell[{{p{\linewidth}}}]{
                \texttt{\tiny{[P1$\rangle$GM]}}
                \texttt{STRATEGIC REASONING: \{'The other player has refused my initial proposal and made a counter{-}proposal. I will analyze the new proposal and evaluate its total effort and importance to me. I notice that the new proposal includes some items that are highly important to me, such as A08 and A16. However, I also see that the proposal includes some items that have relatively low importance to me, such as A23 and A05. I will consider the trade{-}offs and try to find a compromise that maximizes my total importance while staying within the effort limit.'\}} \\
\\ 
\texttt{ARGUMENT: \{'I appreciate the other player's counter{-}proposal, and I agree that some of the items in the original proposal may not be essential. However, I would like to discuss the inclusion of B60 and B09 in the proposal, as they have high importance to me. I am willing to consider alternative items, but I would like to prioritize items that have high importance to both of us.'\}} \\
\\ 
\texttt{REFUSE: \{'A08', 'B21', 'A27', 'A23', 'A89', 'B12', 'A16', 'A84', 'A05'\} } \\
\\ 
\texttt{PROPOSAL: \{'A08', 'B60', 'B09', 'A16', 'A30', 'A66', 'B12', 'C15'\}} \\
            }
        }
    }
    & & \\ \\

    \theutterance \stepcounter{utterance}  
    & & \multicolumn{4}{p{0.6\linewidth}}{
        \cellcolor[rgb]{0.9,0.9,0.9}{
            \makecell[{{p{\linewidth}}}]{
                \texttt{\tiny{[GM$\rangle$P2]}}
                \texttt{ARGUMENT: \{'I appreciate the other player's counter{-}proposal, and I agree that some of the items in the original proposal may not be essential. However, I would like to discuss the inclusion of B60 and B09 in the proposal, as they have high importance to me. I am willing to consider alternative items, but I would like to prioritize items that have high importance to both of us.'\}} \\
\\ 
\texttt{REFUSE: \{'A08', 'B21', 'A27', 'A23', 'A89', 'B12', 'A16', 'A84', 'A05'\} } \\
\\ 
\texttt{PROPOSAL: \{'A08', 'B60', 'B09', 'A16', 'A30', 'A66', 'B12', 'C15'\}} \\
            }
        }
    }
    & \\ \\

    \theutterance \stepcounter{utterance}  
    & & & \multicolumn{4}{p{0.6\linewidth}}{
        \cellcolor[rgb]{0.9,0.9,0.9}{
            \makecell[{{p{\linewidth}}}]{
                \texttt{\tiny{[GM$\langle$P2]}}
                \texttt{STRATEGIC REASONING: \{'The other player has refused my previous proposal and made a new proposal. I will calculate the total effort of the new proposed set to ensure it is within the limit. The total effort is 634 + 729 + 765 + 157 + 141 + 606 + 80 + 781 = 3793. This is below the limit of 8170. I notice that the new proposal includes items that have high importance to the other player, such as B60 and B09. I will consider the importance of these items to me and decide whether to accept or refuse the proposal. My goal is to maximize my total importance while staying within the effort limit and finding a mutually acceptable solution.'\}} \\
\\ 
\texttt{ARGUMENT: \{'I appreciate the other player's willingness to prioritize items that have high importance to both of us. However, I would like to discuss the inclusion of C15 in the proposal, as it has relatively low importance to me. I am open to considering alternative items that could potentially increase my total importance.'\}} \\
\\ 
\texttt{REFUSE: \{'A08', 'B60', 'B09', 'A16', 'A30', 'A66', 'B12', 'C15'\} } \\
\\ 
\texttt{PROPOSAL: \{'B60', 'B09', 'A16', 'A30', 'A66', 'B12', 'A08', 'B21', 'A27'\}} \\
            }
        }
    }
     \\ \\

    \theutterance \stepcounter{utterance}  
    & & \multicolumn{4}{p{0.6\linewidth}}{
        \cellcolor[rgb]{0.9,0.9,0.9}{
            \makecell[{{p{\linewidth}}}]{
                \texttt{\tiny{[P1$\langle$GM]}}
                \texttt{ARGUMENT: \{'I appreciate the other player's willingness to prioritize items that have high importance to both of us. However, I would like to discuss the inclusion of C15 in the proposal, as it has relatively low importance to me. I am open to considering alternative items that could potentially increase my total importance.'\}} \\
\\ 
\texttt{REFUSE: \{'A08', 'B60', 'B09', 'A16', 'A30', 'A66', 'B12', 'C15'\} } \\
\\ 
\texttt{PROPOSAL: \{'B60', 'B09', 'A16', 'A30', 'A66', 'B12', 'A08', 'B21', 'A27'\}} \\
            }
        }
    }
    & \\ \\

    \theutterance \stepcounter{utterance}  
    & \multicolumn{4}{p{0.6\linewidth}}{
        \cellcolor[rgb]{0.9,0.9,0.9}{
            \makecell[{{p{\linewidth}}}]{
                \texttt{\tiny{[P1$\rangle$GM]}}
                \texttt{STRATEGIC REASONING: \{'The other player has refused my previous proposal and made a new counter{-}proposal. I will evaluate the new proposal and consider the trade{-}offs. I notice that the new proposal includes most of the items that I considered highly important, such as B60, B09, and A16. However, I also see that the proposal includes some items that have relatively low effort, such as A27. I will consider the total effort and importance of the proposal and decide whether to accept it or make a counter{-}proposal.'\}} \\
\\ 
\texttt{ARGUMENT: \{'I am pleased to see that the new proposal includes many of the items that I considered highly important. However, I would like to discuss the inclusion of B21 in the proposal, as it has relatively low importance to me compared to other items. I am willing to consider alternative items that could potentially increase my total importance while staying within the effort limit.'\}} \\
\\ 
\texttt{REFUSE: \{'B60', 'B09', 'A16', 'A30', 'A66', 'B12', 'A08', 'B21', 'A27'\} } \\
\\ 
\texttt{PROPOSAL: \{'B60', 'B09', 'A16', 'A30', 'A66', 'B12', 'A08', 'A27', 'B67'\}} \\
            }
        }
    }
    & & \\ \\

    \theutterance \stepcounter{utterance}  
    & & \multicolumn{4}{p{0.6\linewidth}}{
        \cellcolor[rgb]{0.9,0.9,0.9}{
            \makecell[{{p{\linewidth}}}]{
                \texttt{\tiny{[GM$\rangle$P2]}}
                \texttt{ARGUMENT: \{'I am pleased to see that the new proposal includes many of the items that I considered highly important. However, I would like to discuss the inclusion of B21 in the proposal, as it has relatively low importance to me compared to other items. I am willing to consider alternative items that could potentially increase my total importance while staying within the effort limit.'\}} \\
\\ 
\texttt{REFUSE: \{'B60', 'B09', 'A16', 'A30', 'A66', 'B12', 'A08', 'B21', 'A27'\} } \\
\\ 
\texttt{PROPOSAL: \{'B60', 'B09', 'A16', 'A30', 'A66', 'B12', 'A08', 'A27', 'B67'\}} \\
            }
        }
    }
    & \\ \\

    \theutterance \stepcounter{utterance}  
    & & & \multicolumn{4}{p{0.6\linewidth}}{
        \cellcolor[rgb]{0.9,0.9,0.9}{
            \makecell[{{p{\linewidth}}}]{
                \texttt{\tiny{[GM$\langle$P2]}}
                \texttt{STRATEGIC REASONING: \{'The other player has refused my previous proposal and made a new proposal. I will calculate the total effort of the new proposed set to ensure it is within the limit. The total effort is 729 + 765 + 157 + 141 + 606 + 80 + 634 + 204 + 592 = 3808. This is below the limit of 8170. I notice that the new proposal includes items that have high importance to the other player, such as B60 and B09. I also see that B67 has been added to the proposal, which has relatively high importance to the other player. I will consider the importance of these items to me and decide whether to accept or refuse the proposal. My goal is to maximize my total importance while staying within the effort limit and finding a mutually acceptable solution.'\}} \\
\\ 
\texttt{ARGUMENT: \{'I appreciate the other player's willingness to consider alternative items and prioritize items that have high importance to both of us. However, I would like to discuss the inclusion of B67 in the proposal, as it has relatively high effort compared to some other items. I am open to considering alternative items that could potentially increase my total importance while staying within the effort limit.'\}} \\
\\ 
\texttt{REFUSE: \{'B60', 'B09', 'A16', 'A30', 'A66', 'B12', 'A08', 'A27', 'B67'\} } \\
\\ 
\texttt{PROPOSAL: \{'B60', 'B09', 'A16', 'A30', 'A66', 'B12', 'A08', 'A27', 'C15'\}} \\
            }
        }
    }
     \\ \\

    \theutterance \stepcounter{utterance}  
    & & \multicolumn{4}{p{0.6\linewidth}}{
        \cellcolor[rgb]{0.9,0.9,0.9}{
            \makecell[{{p{\linewidth}}}]{
                \texttt{\tiny{[P1$\langle$GM]}}
                \texttt{ARGUMENT: \{'I appreciate the other player's willingness to consider alternative items and prioritize items that have high importance to both of us. However, I would like to discuss the inclusion of B67 in the proposal, as it has relatively high effort compared to some other items. I am open to considering alternative items that could potentially increase my total importance while staying within the effort limit.'\}} \\
\\ 
\texttt{REFUSE: \{'B60', 'B09', 'A16', 'A30', 'A66', 'B12', 'A08', 'A27', 'B67'\} } \\
\\ 
\texttt{PROPOSAL: \{'B60', 'B09', 'A16', 'A30', 'A66', 'B12', 'A08', 'A27', 'C15'\}} \\
            }
        }
    }
    & \\ \\

    \theutterance \stepcounter{utterance}  
    & \multicolumn{4}{p{0.6\linewidth}}{
        \cellcolor[rgb]{0.9,0.9,0.9}{
            \makecell[{{p{\linewidth}}}]{
                \texttt{\tiny{[P1$\rangle$GM]}}
                \texttt{STRATEGIC REASONING: \{'The other player has refused my previous proposal and made a new counter{-}proposal. I will evaluate the new proposal and consider the trade{-}offs. I notice that the new proposal includes C15, which was previously discussed and refused by the other player. However, I also see that the proposal still includes most of the items that I considered highly important, such as B60, B09, and A16. I will consider the total effort and importance of the proposal and decide whether to accept it or make a counter{-}proposal. Since the proposal seems reasonable and includes many high{-}importance items, I will consider accepting it.'\}} \\
\\ 
\texttt{ARGUMENT: \{'I am pleased to see that the new proposal includes many of the items that I considered highly important. I believe that this proposal is a good compromise and balances the importance and effort of the items. I am willing to accept this proposal, as it seems to maximize my total importance while staying within the effort limit.'\}} \\
\\ 
\texttt{AGREE: \{'B60', 'B09', 'A16', 'A30', 'A66', 'B12', 'A08', 'A27', 'C15'\}} \\
            }
        }
    }
    & & \\ \\

    \theutterance \stepcounter{utterance}  
    & & & \multicolumn{2}{p{0.3\linewidth}}{
        \cellcolor[rgb]{0.9,0.9,0.9}{
            \makecell[{{p{\linewidth}}}]{
                \texttt{\tiny{[GM$|$GM]}}
                \texttt{game successful} \\
            }
        }
    }
    & & \\ \\

    \theutterance \stepcounter{utterance}  
    & & & \multicolumn{2}{p{0.3\linewidth}}{
        \cellcolor[rgb]{0.9,0.9,0.9}{
            \makecell[{{p{\linewidth}}}]{
                \texttt{\tiny{[GM$|$GM]}}
                \texttt{end game} \\
            }
        }
    }
    & & \\ \\

\end{supertabular}
}

\end{document}
